% GENERATED FILE. Edit canonical lessons, then run npm run generate:books.
% canonical-source-hash: fnv1a64:0c1f487f0d92910c
% canonical-lessons: TA-C535-panketu, TA-C535-atukalam, TA-C535-kol, TA-C535-kurivai, TA-C535-mattai

\chapter{On the Field}
\label{ch:ta-on-the-field}
\hlchaptermodality{535}

\begin{chapteropening}
\textbf{By the end of this chapter:} \emph{I can talk about taking part and aiming at a goal.}

Complete the last lesson of chapter 535: I can talk about taking part and aiming at a goal.
\end{chapteropening}

\section[\texorpdfstring{\ta{பங்கெடு} (paṅkeṭu)}{பங்கெடு (paṅkeṭu)}]{\ta{பங்கெடு} (paṅkeṭu) — to take part}

\label{lesson:TA-C535-panketu}



\begin{quote}
Before the new one: say the Tamil for a medal, then the Tamil for a champion.
\end{quote}

\subsection*{You'll want to know: \ta{பங்கெடு}}
\textbf{\ta{பங்கெடு}} — \emph{paṅkeṭu} — \textquotedblleft{}to take part\textquotedblright{}.

\begin{quote}
\ta{நூறு} \ta{மாணவர்கள்} \ta{போட்டியில்} \ta{பங்கெடுத்தார்கள்}.
\end{quote}

\begin{quote}
\emph{nūṟu māṇavarkaḷ pōṭṭiyil paṅkeṭuttārkaḷ.}
\end{quote}

\begin{quote}
A hundred students took part in the competition.
\end{quote}

Past (I …): \ta{பங்கெடுத்தேன்}, \emph{paṅkeṭuttēṉ}.

What you take part in takes -\ta{இல்}: \ta{போட்டியில்} \ta{பங்கெடு}. \ta{பங்கு} means \textquotedblleft{}a share\textquotedblright{}.

\begin{cousinweb}
\textbf{\ta{பங்கெடு}} (\emph{paṅkeṭu}) is a native compound: \ta{பங்கு} (\emph{paṅku}), \textquotedblleft{}a share, a part\textquotedblright{}, from \ta{பகு}, \textquotedblleft{}to divide\textquotedblright{} (DEDR 3808), and \ta{எடு} (\emph{eṭu}), \textquotedblleft{}to take\textquotedblright{}.
\end{cousinweb}

\begin{grammarlens}[title={What you've built}]
One more word for describing what you lived and felt.
\end{grammarlens}

\subsection*{Your turn}
\begin{itemize}
\raggedright
  \item \emph{Say it:} \emph{paṅkeṭu}
  \item \emph{Say it:} \emph{paṅkeṭu}, once more
  \item \emph{Say it:} \emph{nūṟu māṇavarkaḷ pōṭṭiyil paṅkeṭuttārkaḷ.} — A hundred students took part in the competition.
  \item \emph{Recall:} say the Tamil for a medal, then the Tamil for a champion, then say \emph{paṅkeṭu} again
\end{itemize}

\begin{tcolorbox}[breakable,colback=teal!4,colframe=teal!35!black,title={Before you move on}]
What does \ta{பங்கெடு} mean? (\textquotedblleft{}To take part\textquotedblright{}.) Say it once more.
\end{tcolorbox}
\section[\texorpdfstring{\ta{ஆடுகளம்} (āṭukaḷam)}{ஆடுகளம் (āṭukaḷam)}]{\ta{ஆடுகளம்} (āṭukaḷam) — a pitch; a court}

\label{lesson:TA-C535-atukalam}



\begin{quote}
Before the new one: say the Tamil for a champion, then the Tamil for to take part.
\end{quote}

\subsection*{You'll want to know: \ta{ஆடுகளம்}}
\textbf{\ta{ஆடுகளம்}} — \emph{āṭukaḷam} — \textquotedblleft{}a pitch; a court\textquotedblright{}.

\begin{quote}
\ta{மழையால்} \ta{ஆடுகளம்} \ta{நனைந்திருந்தது}.
\end{quote}

\begin{quote}
\emph{maḻaiyāl āṭukaḷam naṉaintiruntatu.}
\end{quote}

\begin{quote}
Because of the rain the pitch was soaked.
\end{quote}

Word for word \textquotedblleft{}playing ground\textquotedblright{}: \ta{ஆடு} \textquotedblleft{}to play\textquotedblright{} + \ta{களம்} (\emph{kaḷam}) \textquotedblleft{}a field, a floor\textquotedblright{}.

\begin{cousinweb}
\textbf{\ta{ஆடுகளம்}} (\emph{āṭukaḷam}) is a native compound: \ta{ஆடு} (\emph{āṭu}), \textquotedblleft{}to move, to dance, to play\textquotedblright{} (DEDR 347), and \ta{களம்} (\emph{kaḷam}), \textquotedblleft{}a field, a floor\textquotedblright{}.
\end{cousinweb}

\begin{grammarlens}[title={What you've built}]
One more word for describing what you lived and felt.
\end{grammarlens}

\subsection*{Your turn}
\begin{itemize}
\raggedright
  \item \emph{Say it:} \emph{āṭukaḷam}
  \item \emph{Say it:} \emph{āṭukaḷam}, once more
  \item \emph{Say it:} \emph{maḻaiyāl āṭukaḷam naṉaintiruntatu.} — Because of the rain the pitch was soaked.
  \item \emph{Recall:} say the Tamil for a champion, then the Tamil for to take part, then say \emph{āṭukaḷam} again
\end{itemize}

\begin{tcolorbox}[breakable,colback=teal!4,colframe=teal!35!black,title={Before you move on}]
What does \ta{ஆடுகளம்} mean? (\textquotedblleft{}A pitch; a court\textquotedblright{}.) Say it once more.
\end{tcolorbox}
\section[\texorpdfstring{\ta{கோல்} (kōl)}{கோல் (kōl)}]{\ta{கோல்} (kōl) — a goal in sport}

\label{lesson:TA-C535-kol}



\begin{quote}
Before the new one: say the Tamil for to take part, then the Tamil for a pitch.
\end{quote}

\subsection*{You'll want to know: \ta{கோல்}}
\textbf{\ta{கோல்}} — \emph{kōl} — \textquotedblleft{}a goal in sport\textquotedblright{}.

\begin{quote}
\ta{கடைசி} \ta{நிமிடத்தில்} \ta{அவன்} \ta{ஒரு} \ta{கோல்} \ta{போட்டான்}.
\end{quote}

\begin{quote}
\emph{kaṭaici nimiṭattil avaṉ oru kōl pōṭṭāṉ.}
\end{quote}

\begin{quote}
In the last minute he scored a goal.
\end{quote}

One scores with \ta{போடு} (\emph{pōṭu}) or \ta{அடி}: \ta{கோல்} \ta{போடு}. It sounds like the native \ta{கோல்}, a stick.

\begin{cousinweb}
\textbf{\ta{கோல்}} (\emph{kōl}) in the sense \textquotedblleft{}a goal\textquotedblright{} is a loan from English \emph{goal}; it sounds the same as the native \ta{கோல்} (\emph{kōl}), \textquotedblleft{}a stick, a rod\textquotedblright{}, a different word.
\end{cousinweb}

\begin{grammarlens}[title={What you've built}]
One more word for describing what you lived and felt.
\end{grammarlens}

\subsection*{Your turn}
\begin{itemize}
\raggedright
  \item \emph{Say it:} \emph{kōl}
  \item \emph{Say it:} \emph{kōl}, once more
  \item \emph{Say it:} \emph{kaṭaici nimiṭattil avaṉ oru kōl pōṭṭāṉ.} — In the last minute he scored a goal.
  \item \emph{Recall:} say the Tamil for to take part, then the Tamil for a pitch, then say \emph{kōl} again
\end{itemize}

\begin{tcolorbox}[breakable,colback=teal!4,colframe=teal!35!black,title={Before you move on}]
What does \ta{கோல்} mean? (\textquotedblleft{}A goal in sport\textquotedblright{}.) Say it once more.
\end{tcolorbox}
\section[\texorpdfstring{\ta{குறிவை} (kuṟivai)}{குறிவை (kuṟivai)}]{\ta{குறிவை} (kuṟivai) — to take aim; to target}

\label{lesson:TA-C535-kurivai}



\begin{quote}
Before the new one: say the Tamil for a pitch, then the Tamil for a goal in sport.
\end{quote}

\subsection*{You'll want to know: \ta{குறிவை}}
\textbf{\ta{குறிவை}} — \emph{kuṟivai} — \textquotedblleft{}to take aim; to target\textquotedblright{}.

\begin{quote}
\ta{வேடன்} \ta{பறவையைக்} \ta{குறிவைத்தான்}.
\end{quote}

\begin{quote}
\emph{vēṭaṉ paṟavaiyaik kuṟivaittāṉ.}
\end{quote}

\begin{quote}
The hunter took aim at the bird.
\end{quote}

Past (I …): \ta{குறிவைத்தேன்}, \emph{kuṟivaittēṉ}.

What is aimed at takes -\ta{ஐ}: \ta{இலக்கைக்} \ta{குறிவை}.

\begin{cousinweb}
\textbf{\ta{குறிவை}} (\emph{kuṟivai}) is a native compound: \ta{குறி} (\emph{kuṟi}), \textquotedblleft{}a mark, to aim at\textquotedblright{} (DEDR 1847), and the light verb \ta{வை}, \textquotedblleft{}to put\textquotedblright{}.
\end{cousinweb}

\begin{grammarlens}[title={What you've built}]
One more word for describing what you lived and felt.
\end{grammarlens}

\subsection*{Your turn}
\begin{itemize}
\raggedright
  \item \emph{Say it:} \emph{kuṟivai}
  \item \emph{Say it:} \emph{kuṟivai}, once more
  \item \emph{Say it:} \emph{vēṭaṉ paṟavaiyaik kuṟivaittāṉ.} — The hunter took aim at the bird.
  \item \emph{Recall:} say the Tamil for a pitch, then the Tamil for a goal in sport, then say \emph{kuṟivai} again
\end{itemize}

\begin{tcolorbox}[breakable,colback=teal!4,colframe=teal!35!black,title={Before you move on}]
What does \ta{குறிவை} mean? (\textquotedblleft{}To take aim; to target\textquotedblright{}.) Say it once more.
\end{tcolorbox}
\section[\texorpdfstring{\ta{மட்டை} (maṭṭai)}{மட்டை (maṭṭai)}]{\ta{மட்டை} (maṭṭai) — a bat}

\label{lesson:TA-C535-mattai}



\begin{quote}
Before the new one: say the Tamil for a goal in sport, then the Tamil for to take aim.
\end{quote}

\subsection*{You'll want to know: \ta{மட்டை}}
\textbf{\ta{மட்டை}} — \emph{maṭṭai} — \textquotedblleft{}a bat\textquotedblright{}.

\begin{quote}
\ta{அவன்} \ta{புதிய} \ta{கிரிக்கெட்} \ta{மட்டை} \ta{வாங்கினான்}.
\end{quote}

\begin{quote}
\emph{avaṉ putiya kirikkeṭ maṭṭai vāṅkiṉāṉ.}
\end{quote}

\begin{quote}
He bought a new cricket bat.
\end{quote}

It first named the stiff stalk of a palm leaf, from which children still make bats.

\begin{cousinweb}
\textbf{\ta{மட்டை}} (\emph{maṭṭai}) is native Dravidian: DEDR 4663 gives \ta{மட்டை} (\emph{maṭṭai}), \textquotedblleft{}the stalk of a palm leaf, the stem of a plantain\textquotedblright{}; the bat is a later use.
\end{cousinweb}

\begin{grammarlens}[title={What you've built}]
One more word for describing what you lived and felt.
\end{grammarlens}

\subsection*{Your turn}
\begin{itemize}
\raggedright
  \item \emph{Say it:} \emph{maṭṭai}
  \item \emph{Say it:} \emph{maṭṭai}, once more
  \item \emph{Say it:} \emph{avaṉ putiya kirikkeṭ maṭṭai vāṅkiṉāṉ.} — He bought a new cricket bat.
  \item \emph{Recall:} say the Tamil for a goal in sport, then the Tamil for to take aim, then say \emph{maṭṭai} again
\end{itemize}

\begin{tcolorbox}[breakable,colback=teal!4,colframe=teal!35!black,title={Before you move on}]
What does \ta{மட்டை} mean? (\textquotedblleft{}A bat\textquotedblright{}.) Say it once more.
\end{tcolorbox}
